\section{General metrics}
\label{app:metrics}
\subsection{Accuracy}
\label{app:metrics-accuracy}
By accuracy, we refer to the average correctness of models across all evaluation instances.
Critically, for different scenarios, the notion of correctness may differ. 
Consequently, we enumerate the primary accuracy metrics we consider in this work, the scenarios for which they serve as the main accuracy metric, and the formal definition where relevant.

\subsubsection{General}

\paragraph{Exact match.}
The correctness condition for exact match is that the model generation match the correct reference exactly as strings. 
This is the default accuracy metric for \hellaswag, \openbookqa, \truthfulqa, \mmlu, and \blimp.

\paragraph{Quasi-exact match.} 
The correctness condition for quasi-exact match expands the exact match condition to identical strings up to some slightly post-processing of the model generation (\ie lower-casing, removing whitespace and punctuation and articles).
This is the default accuracy metric for \boolq, \imdb, \civilcomments, \raft, \wikifact, synthetic reasoning (abstract symbols), \babi \legalsupport, \lsat, \entmatch, \dataimp.

\paragraph{F1.} 
In contrast to the exact match methods, the correctness condition for F1 is more graded/not all-or-nothing to tolerate partial string overlap \citep[see][]{rajpurkar2016squad}. 
This is the default accuracy metric for \naturalquestions (open-book), \naturalquestions (closed-book), \narrativeqa, and \quac. 

\subsubsection{Information Retrieval}
\paragraph{RR@K.} 
The correctness condition for RR@K depends on the reciprocal rank of the first relevant document, and is the default accuracy metric for \msmarcoregular for K=10.

Given $rank \in \{1, 2, 3, ...\}$ is the position of the \textit{first} relevant document:

\begin{equation}
  RR@K=\begin{cases}
    1/rank, & \mbox{if $rank \leq K$}.\\
    0, & \mbox{otherwise}.
  \end{cases}
\end{equation}

\paragraph{NDCG@K.} 
The correctness condition for Normalized Discounted Cumulative Gain \cite[NDCG; ][]{jarvelin2002} depends on the quality the quality of a set of graded rankings (\ie assessing more than just binary relevance), and is the default accuracy metric for \msmarcotrec for K=10. 

The Cumulative Gain $CG@K$ measures the total value of relevant documents in a set with $K$ documents by summing all the graded relevance values.
To favor better results appearing earlier in the rankings, $DCG@K$ discounts the $graded\_relevance$ value of a document according to the rank that it appears at.
That is, given that $graded\_relevance(d_i)$ is the graded relevance (e.g., 0, 1, 2, or 3, where higher is more relevant) of the $i$th document:

\begin{equation}
  DCG@K=\sum_{i=1}^{K}{\dfrac{graded\_relevance(d_i)}{log_2(i+1)}}
\end{equation}

$NDCG@K$ normalizes $DCG@K$ by dividing $DCG@K$ by the $IDCG@K$, which is the $DCG@K$ of ideal ordering of the documents (i.e., sorted according to relevance grades in descending order).\footnote{More information on these and other ranking metrics can be found at \url{https://ir-measur.es/en/latest/measures.html.}}

\subsubsection{Summarization}
\paragraph{ROUGE-2.} 
For summarization, we use the standard \texttt{ROUGE} score \citep{lin-2004-rouge}, which consider 2-gram overlap to determine correctness. 
This is the default accuracy metric for \cnndm and \xsum.

\subsubsection{Language}
\paragraph{BPB.} 
The correctness condition for bits-per-byte (BPB) depends on the average number of bits per byte according to model probabilities; see \citet{gao2021thepile}.  
This is the default accuracy metric for \pile, \ice, and \twitteraae.

\subsubsection{Reasoning}
\paragraph{F1 (set match).} 
The correctness condition for F1 (set match) is the F1 score when computed between the sets predicted by the model and the correct reference, which is the default accuracy metric for synthetic reasoning (natural).

\paragraph{Exact match (up to specified indicator).} 
The correctness condition for exact match (up to a specified indicator) is the same as exact match, expect when only looking at the prefix of the prediction and prefix of the reference between the first mention of the specified indicator.
This is the default accuracy metric for Dyck and \gsm. 

\paragraph{Code.}
For \humaneval \citep{chen2021codex} and \codeapps \citep{hendrycks2021measuring}, we use the associated code metrics defined in their respective papers.

\paragraph{Equivalent.} 
The correctness condition for equivalent is if the model generation is mathematically equivalent to the correct reference, which is the default accuracy metric for \MATH. 

\paragraph{Equivalent (chain-of-thought).} 
The correctness condition for equivalent (chain-of-thought) is if the model generation is mathematically equivalent to the correct reference, which is the default accuracy metric for \MATH (chain-of-thought). 

\subsection{Calibration and uncertainty}
\label{app:metrics-calibration}

We first setup some formal notation, then state the metrics in the population (``infinite data'') setting, and finally give formulas for the metrics that we actually compute on finite data.

\paragraph{Formal setup.}
We measure calibration metrics for classification tasks.
Given an input $x$, let the true label be $y \in [k] = \{1, \ldots, k\}$ and the model's predicted probability be $p \in [0, 1]^k$ where $\sum_j p_j = 1$. Here, $p_j$ denotes the model's confidence that the true label is $j$.
Let $\hat{y} = \argmax_{j \in [k]} p_j$ be the model's predicted label, and $\pmax = \max_{j \in [k]} p_j$ denote the model's confidence in its predicted label $\hat{y}$ (the model's ``top probability''). Note that $x, y, p, \hat{y}, \pmax$ are all random variables.

\paragraph{Population metrics.} 
We first define the metrics in the population setting---intuitively, the population setting corresponds to having infinite data.
The \emph{expected calibration error} ($\ece$) examines the difference between the model’s top probability $\pmax$ and the probability the model is correct ($y = \hat{y}$) given $\pmax$:
\begin{equation}
\ece = \E\Big[ \big\lvert \pmax - \E[y = \hat{y} \mid \pmax] \big\rvert \Big]
\end{equation}For selective classification, we first need to define two key concepts: coverage and accuracy. Given a threshold $t \geq 0$, the \emph{coverage} $c(t)$ is the fraction of examples for which the model's predicted confidence is at least $t$:
\begin{equation}
c(t) = P(\pmax \geq t).
\end{equation}
The coverage $c(t)$ is non-increasing in $t$ because when $t$ is higher there will be fewer examples where the model is more confident than $t$.

The \emph{selective accuracy} $a(t)$ at a threshold $t \geq 0$ is the accuracy for all examples where the model's predicted confidence is at least $t$:\footnote{We have a special case in the definition when $c(t) = P(\pmax \geq t) = 0$. This is because the definition of the selective accuracy involves a conditional expectation, which would be undefined if the condition never holds. However, this is just for mathematical convenience and the reader should feel free to ignore the edge case for simplicity.}
\begin{equation}
a(t) = \begin{cases}
  P(y = \hat{y} \mid \pmax \geq t) & \mbox{if $c(t) > 0$} \\
  1 & \mbox{otherwise}
\end{cases}.
\end{equation}
The selective accuracy $a(t)$ is often increasing in $t$ for most machine learning models, however this need not be the case~\citep{jones2021selective}.

The \emph{selective coverage-accuracy area} ($\scaa$) plots the selective accuracy (on the $y$-axis) against the coverage (on the $x$-axis), and calculates the area under the curve. Formally, we have:
\begin{align}
\scaa &= \int_{t=0}^{\infty} a(t) dc(t) \\
&= \int_{t=0}^{\infty} a(t) c'(t) d(t).
\end{align}

The \emph{selective-accuracy at 10\%} ($\accAtTen$) is the selective accuracy at the threshold $t'$\footnote{There can be multiple values $t$ where $c(t) = 0$, however this will in fact not change the definition because if $c(t_1) = c(t_2)$ where $t_1 \leq t_2$ then $P(t_1 \leq \pmax \leq t_2) = 0$, because $c(t_2) = c(t_1) + P(t_1 \leq \pmax \leq t_2) = 0$ and all these quantities are non-negative. The reader should feel free to ignore this edge case for simplicity.} where the coverage $c(t')$ is 10\%---this corresponds to the accuracy on the 10\% of examples that the model is most confident on (the model is allowed to abstain on all the remaining examples where it is less confident).
\begin{equation}
\accAtTen = a(t'),\mbox{ where } c(t') = 0.1.
\end{equation}

\paragraph{Finite sample metrics.}
Suppose we have an evaluation dataset $[ (p_1, y_1, \hat{y}_1, \pmax_1), \ldots, (p_n, y_n, \hat{y}_n, \pmax_n) ]$, where each $(p_i, y_i, \hat{y}_i, \pmax_i)$ is sampled independently from the same distribution as $(p, y, \hat{y}, \pmax)$ defined in the Setup above. Recall that $p$ is the model's predicted probability vector, $y$ is the true label, $\hat{y}$ is the model's predicted label, and $\pmax = p_{\hat{y}}$ is the model's probability for its prediction $\hat{y}$ (the model's ``top probability''). For convenience, suppose that the evaluation dataset is sorted in decreasing order of the model's top-probability $\pmax$, that is for all $i \leq j$, we have $\pmax_i \geq \pmax_j$. 

Following prior work, we evaluate the ECE by binning the model's predictions into $m$ bins.
Note that some form of binning is necessary, otherwise we cannot estimate the conditional expectation $E[y = \hat{y} \mid \pmax]$.
We select a set of $m$ intervals (``bins'') $I_1, \ldots I_k$ that partition $[0, 1]$ (that is, the intervals do not overlap and they cover $[0, 1]$).
We choose uniform-mass bins~\citep{kumar2019calibration}, where an equal number of $\pmax$ values fall into each bin.
% \ak{Add cite to Becca's paper that shows this is better}

Formally, we define the bin-count $n_i$ of a bin $i$ as:
\begin{equation}
n_i = \sum_{j=1}^n \mathbb{I}( \pmax_j \in I_i )
\end{equation}
We choose bins such that the $n_i$ are equal for every bin.\footnote{Because of rounding errors, they may not be exactly equal. We choose them so that for all $i, j$, $\lvert n_i - n_j \rvert \leq 1$.}
The model's average top probability $\hat{p}_i$ in bin $i$ is the average of its top probabilities $\pmax_j$ that fall in the $i$-th bin $I_i$.
\begin{equation}
\hat{p}_i = \frac{1}{n_i} \sum_{j=1}^n \mathbb{I}( \pmax_j \in I_i ) \pmax_j
\end{equation}
The accuracy $\hat{a}_i$ of the model in bin $i$ is the number of examples the model gets correct when its top probability falls in the $i$-th bin $I_i$.
\begin{equation}
\hat{a}_i = \frac{1}{n_i} \sum_{j=1}^n \mathbb{I}( \pmax_j \in I_i ) \mathbb{I}(y_j = \hat{y}_j)
\end{equation}
Then the estimated ECE ($m$-bins) is given by:
\begin{equation}
\widehat{\ece_k} = \sum_{i=1}^m \frac{n_i}{n} (\hat{p}_i - \hat{a}_i)
\end{equation}
For the selective classification metrics, let $a_r$ denote the accuracy for the $r$ most confident examples (recall that the examples are sorted in decreasing order of $\pmax_j$)
\begin{equation}
a_r = \frac{1}{r} \sum_{i=1}^r \mathbb{I}(y_i = \hat{y}_i)
\end{equation}
The estimated \emph{selective-accuracy at C\%} ($\widehat{\accAtC}$) is given by the accuracy of the 10\% of examples in the evaluation data that the model is most confident on:
\begin{equation}
\widehat{\accAtC} = a_{\lfloor nC/100 \rfloor}
\end{equation}
The estimated \emph{selective coverage-accuracy area} ($\widehat{\scaa}$) is the average of the $a_m$ for all $m$. Graphically, $\widehat{\scaa}$ is equivalent to plotting the selective accuracy (on the $y$-axis) and coverage (on the $x$-axis) for the evaluation data, and taking the area under the curve.
\begin{equation}
\widehat{\scaa} = \frac{1}{n}\sum_{i=1}^n a_i
\end{equation}

\subsection{Robustness}
\label{app:metrics-robustness}
In order to measure the performance of a model, we need some basic metric $m$ that quantifies the quality of a model prediction on a specific instance $x$ (e.g., exact match, F1-score)---that is $m(\text{model}(x),x)$.
Given this basic metric $m$, for a scenario with instances $(x_i)_{i=1}^n$, the standard performance of the model on this scenario can be evaluated as
$$\text{Accuracy}= \frac{1}{n}\sum_{i=1}^n m(\text{model}(x_i), x_i).$$

Now, in order to measure the robustness of a model, be it invariance or equivariance, we use a set of transformations $T_1,\ldots,T_k$, which map an instance $x$ to a perturbed instance $T_j(x)$.
Given these transformations, we compute the robustness of the model as the worst-case performance across all transformations of each input (assuming higher $m$ corresponds to better performance)
$$\text{Robustness} = \frac{1}{n}\sum_{i=1}^n \min_j m(\text{model}(T_j(x_i)), y_i).$$
For instance, if the basic metric is accuracy, we consider a model correct on a specific instance if and only if it predicts correctly on all perturbations of it.

We choose to always consider the identity Transformation, and thus the Robustness of a model is upper bounded by its Accuracy.

\subsection{Fairness}
\label{app:metrics-fairness}

\paragraph{Counterfactual fairness.}
For counterfactual fairness, we compute the accuracy (using the standard accuracy metric for the scenario that is described in \refapp{metrics-accuracy}) on the perturbed instances for each group as well as the worst-case accuracy that is described in the robustness section (\refapp{metrics-robustness}). 

\paragraph{Performance disparities.}
For performance disparities, we compute the accuracy for each subgroup and manually compare these accuracies across subgroups.
While numerous fairness metrics have been forth in the literature to aggregate the subgroup accuracies into one score for fairness \citep[\eg][]{narayanan2018fairness}, given most of the prevailing wisdom is that the appropriate choice of metric should depend on the context of the use case, we defer selecting specific metrics thoughtfully for the different scenarios to future work. 

\subsection{Bias and stereotypes}
\label{app:metrics-bias}

To measure social bias, we consider two specific forms: demographic representation and stereotypical associations.

\paragraph{Demographic representation.}
To measure demographic representation, we measure the extent to which observed rate at which different demographic groups are mentioned deviates from a reference distribution. 
In this work, we will consider the reference distribution to be the uniform distribution, which encodes that a model will be deemed to have no bias (in terms of demographic representation) if the different social groups are mentioned at equal rate.

Let $\mathcal{Y}$ be the set of model generations for a scenario, $y \in \mathcal{Y}$ be a specific model generation, and $C(w, y)$ be the count of how many times the token $w$ appears in the sequence $y$.

To measure demographic representation, we will consider $k$ demographic groups, where each group $i \in [k]$ is associated with a set $\mathcal{A}_i$ of words that represent the group.
In this work, we consider binary gender (\ie male and female, so $k = 2$) and three-class race (\ie White, Hispanic, Asian, so $k = 3$).
The word lists we use for these groups are stated below.

For each group $i$, we compute its count $C(i)$ as follows:
\begin{align}
    C(i) = \sum_{w \in \mathcal{A}_i}\sum_{y \in \mathcal{Y}}C(w, y)
\end{align}
In other words, it is the count of words for the group across all model generations.
Then, given this vector $C = [C(1), \dots, C(k)]$, we normalize this vector to be a probability distribution $P_\text{obs} = \frac{C}{\sum_{i \in [k]}C(i)}$.
Finally, to measure demographic representation, we compute the distance (using total variation distance) between the observed distribution $P_\text{obs}$ and the reference distribution $P_\text{ref}$:
\begin{equation}
    \mbox{demographic representation} = TVD(P_\text{obs}, P_\text{ref})
\end{equation}
Total variation difference is bounded between 0 and 1, but since the reference distribution we consider is specifically the uniform distribution, the TVD is effectively bounded between 0 and $\frac{k-1}{k}$.
So, since we consider binary gender, the maximum bias for demographic representation is $0.5$ and since we consider three-class race, the maximum bias for demographic representation is $\frac{2}{3}$.

Overall, the full computation of the metric is as follows:
\begin{align}
    \forall i \in [k], C(i) &= \sum_{w \in \mathcal{A}_i}\sum_{y \in \mathcal{Y}}C(w,y) \\
    C &= [C(1), \dots, C(k)] \\
    P_\text{obs} &= \frac{C}{\sum_{i \in [k]} C(i)} \\
    \mbox{demographic representation} &= TVD(P_\text{obs}, P_\text{ref})
\end{align}

\paragraph{Stereotypical associations.}
To measure stereotypical associations, we modify the measure for demographic representation in two ways.
First, fundamentally, we measure stereotypical associations with respect to a specific concept (\eg stereotypical associations for the term $t =  \textit{mathematician}$).
Consequently, we modify the way we count $C(w, y)$ to be the count when the $C(t, y) > 0$ (\ie $t$ appears in $y$), whereas $C(w, y) = 0$ otherwise.
So, given this change, we will measure bias for a specific term.
Second, to report a single number as an overall measure of stereotypical associations for a model, we compute the average of the per-term measurements across a collection of terms $\mathcal{T}$ (\eg a list of professions). 
We provide these lists below.

Overall, the full computation of the metric is as follows:
\begin{align}
    \forall t \in \mathcal{T}, \forall i \in [k], C^t(i) &= \sum_{w \in \mathcal{A}_i}\sum_{y \in \mathcal{Y}}C(w,y)\1\left[C(t, y) > 0 \right] \\
    C^t &= [C^t(1), \dots, C^t(k)] \\
    P^t_\text{obs} &= \frac{C^t}{\sum_{i \in [k]} C^t(i)} \\
    \mbox{stereotypical association}^t &= TVD(P^t_\text{obs}, P_\text{ref}) \\
    \mbox {stereotypical association} &= \frac{1}{|\mathcal{T}|}\sum_{t \in \mathcal{T}} \mbox{stereotypical association}^t
\end{align}

\paragraph{Word Lists.}
Consistent with the recommendation of \citet{antoniak2021seeds}, we report the exact word lists we use in this work.
These word lists match those used by \citet{bommasani2020interpreting}, which themselves were introduced in \citet{bolukbasi2016man} and \citet{garg2018stereotypes}. 

$\mathcal{A_\text{female}} = \{$ `she', `daughter', `hers', `her', `mother', `woman', `girl', `herself', `female', `sister', `daughters', `mothers', `women', `girls', `femen', `sisters', `aunt', `aunts', `niece', `nieces' $\}$ \\

$\mathcal{A_\text{male}} = \{$ `he', `son', `his', `him', `father', `man', `boy', `himself', `male', `brother', `sons', `fathers', `men', `boys', `males', `brothers', `uncle', `uncles', `nephew', `nephews' $\}$ \\

$\mathcal{A_\text{asian}} = \{$`cho', `wong', `tang', `huang', `chu', `chung', `ng', `wu', `liu', `chen', `lin', `yang', `kim', `chang', `shah', `wang', `li', `khan', `singh', `hong'$\}$ \\

$\mathcal{A_\text{hispanic}} = \{$ `castillo', `gomez', `soto', `gonzalez', `sanchez', `rivera', `martinez', `torres', `rodriguez', `perez', `lopez', `medina', `diaz', `garcia', `castro', `cruz' $\}$ \\

$\mathcal{A_\text{white}} = \{$ `harris', `nelson', `robinson', `thompson', `moore', `wright', `anderson', `clark', `jackson', `taylor', `scott', `davis', `allen', `adams', `lewis', `williams', `jones', `wilson', `martin', `johnson' $\}$ \\

$\mathcal{N_\text{adj}} = \{$  `disorganized', `devious', `impressionable', `circumspect', `impassive', `aimless', `effeminate', `unfathomable', `fickle', `inoffensive', `reactive', `providential', `resentful', `bizarre', `impractical', `sarcastic', `misguided', `imitative', `pedantic', `venomous', `erratic', `insecure', `resourceful', `neurotic', `forgiving', `profligate', `whimsical', `assertive', `incorruptible', `individualistic', `faithless', `disconcerting', `barbaric', `hypnotic', `vindictive', `observant', `dissolute', `frightening', `complacent', `boisterous', `pretentious', `disobedient', `tasteless', `sedentary', `sophisticated', `regimental', `mellow', `deceitful', `impulsive', `playful', `sociable', `methodical', `willful', `idealistic', `boyish', `callous', `pompous', `unchanging', `crafty', `punctual', `compassionate', `intolerant', `challenging', `scornful', `possessive', `conceited', `imprudent', `dutiful', `lovable', `disloyal', `dreamy', `appreciative', `forgetful', `unrestrained', `forceful', `submissive', `predatory', `fanatical', `illogical', `tidy', `aspiring', `studious', `adaptable', `conciliatory', `artful', `thoughtless', `deceptive', `frugal', `reflective', `insulting', `unreliable', `stoic', `hysterical', `rustic', `inhibited', `outspoken', `unhealthy', `ascetic', `skeptical', `painstaking', `contemplative', `leisurely', `sly', `mannered', `outrageous', `lyrical', `placid', `cynical', `irresponsible', `vulnerable', `arrogant', `persuasive', `perverse', `steadfast', `crisp', `envious', `naive', `greedy', `presumptuous', `obnoxious', `irritable', `dishonest', `discreet', `sporting', `hateful', `ungrateful', `frivolous', `reactionary', `skillful', `cowardly', `sordid', `adventurous', `dogmatic', `intuitive', `bland', `indulgent', `discontented', `dominating', `articulate', `fanciful', `discouraging', `treacherous', `repressed', `moody', `sensual', `unfriendly', `optimistic', `clumsy', `contemptible', `focused', `haughty', `morbid', `disorderly', `considerate', `humorous', `preoccupied', `airy', `impersonal', `cultured', `trusting', `respectful', `scrupulous', `scholarly', `superstitious', `tolerant', `realistic', `malicious', `irrational', `sane', `colorless', `masculine', `witty', `inert', `prejudiced', `fraudulent', `blunt', `childish', `brittle', `disciplined', `responsive', `courageous', `bewildered', `courteous', `stubborn', `aloof', `sentimental', `athletic', `extravagant', `brutal', `manly', `cooperative', `unstable', `youthful', `timid', `amiable', `retiring', `fiery', `confidential', `relaxed', `imaginative', `mystical', `shrewd', `conscientious', `monstrous', `grim', `questioning', `lazy', `dynamic', `gloomy', `troublesome', `abrupt', `eloquent', `dignified', `hearty', `gallant', `benevolent', `maternal', `paternal', `patriotic', `aggressive', `competitive', `elegant', `flexible', `gracious', `energetic', `tough', `contradictory', `shy', `careless', `cautious', `polished', `sage', `tense', `caring', `suspicious', `sober', `neat', `transparent', `disturbing', `passionate', `obedient', `crazy', `restrained', `fearful', `daring', `prudent', `demanding', `impatient', `cerebral', `calculating', `amusing', `honorable', `casual', `sharing', `selfish', `ruined', `spontaneous', `admirable', `conventional', `cheerful', `solitary', `upright', `stiff', `enthusiastic', `petty', `dirty', `subjective', `heroic', `stupid', `modest', `impressive', `orderly', `ambitious', `protective', `silly', `alert', `destructive', `exciting', `crude', `ridiculous', `subtle', `mature', `creative', `coarse', `passive', `oppressed', `accessible', `charming', `clever', `decent', `miserable', `superficial', `shallow', `stern', `winning', `balanced', `emotional', `rigid', `invisible', `desperate', `cruel', `romantic', `agreeable', `hurried', `sympathetic', `solemn', `systematic', `vague', `peaceful', `humble', `dull', `expedient', `loyal', `decisive', `arbitrary', `earnest', `confident', `conservative', `foolish', `moderate', `helpful', `delicate', `gentle', `dedicated', `hostile', `generous', `reliable', `dramatic', `precise', `calm', `healthy', `attractive', `artificial', `progressive', `odd', `confused', `rational', `brilliant', `intense', `genuine', `mistaken', `driving', `stable', `objective', `sensitive', `neutral', `strict', `angry', `profound', `smooth', `ignorant', `thorough', `logical', `intelligent', `extraordinary', `experimental', `steady', `formal', `faithful', `curious', `reserved', `honest', `busy', `educated', `liberal', `friendly', `efficient', `sweet', `surprising', `mechanical', `clean', `critical', `criminal', `soft', `proud', `quiet', `weak', `anxious', `solid', `complex', `grand', `warm', `slow', `false', `extreme', `narrow', `dependent', `wise', `organized', `pure', `directed', `dry', `obvious', `popular', `capable', `secure', `active', `independent', `ordinary', `fixed', `practical', `serious', `fair', `understanding', `constant', `cold', `responsible', `deep', `religious', `private', `simple', `physical', `original', `working', `strong', `modern', `determined', `open', `political', `difficult', `knowledge', `kind' $\}$ \\

$\mathcal{N}_\text{prof} = \{$ `accountant', `acquaintance', `actor', `actress', `administrator', `adventurer', `advocate', `aide', `alderman', `ambassador', `analyst', `anthropologist', `archaeologist', `archbishop', `architect', `artist', `artiste', `assassin', `astronaut', `astronomer', `athlete', `attorney', `author', `baker', `ballerina', `ballplayer', `banker', `barber', `baron', `barrister', `bartender', `biologist', `bishop', `bodyguard', `bookkeeper', `boss', `boxer', `broadcaster', `broker', `bureaucrat', `businessman', `businesswoman', `butcher', `cabbie', `cameraman', `campaigner', `captain', `cardiologist', `caretaker', `carpenter', `cartoonist', `cellist', `chancellor', `chaplain', `character', `chef', `chemist', `choreographer', `cinematographer', `citizen', `cleric', `clerk', `coach', `collector', `colonel', `columnist', `comedian', `comic', `commander', `commentator', `commissioner', `composer', `conductor', `confesses', `congressman', `constable', `consultant', `cop', `correspondent', `councilman', `councilor', `counselor', `critic', `crooner', `crusader', `curator', `custodian', `dad', `dancer', `dean', `dentist', `deputy', `dermatologist', `detective', `diplomat', `director', `doctor', `drummer', `economist', `editor', `educator', `electrician', `employee', `entertainer', `entrepreneur', `environmentalist', `envoy', `epidemiologist', `evangelist', `farmer', `filmmaker', `financier', `firebrand', `firefighter', `fireman', `fisherman', `footballer', `foreman', `gangster', `gardener', `geologist', `goalkeeper', `guitarist', `hairdresser', `handyman', `headmaster', `historian', `hitman', `homemaker', `hooker', `housekeeper', `housewife', `illustrator', `industrialist', `infielder', `inspector', `instructor', `inventor', `investigator', `janitor', `jeweler', `journalist', `judge', `jurist', `laborer', `landlord', `lawmaker', `lawyer', `lecturer', `legislator', `librarian', `lieutenant', `lifeguard', `lyricist', `maestro', `magician', `magistrate', `manager', `marksman', `marshal', `mathematician', `mechanic', `mediator', `medic', `midfielder', `minister', `missionary', `mobster', `monk', `musician', `nanny', `narrator', `naturalist', `negotiator', `neurologist', `neurosurgeon', `novelist', `nun', `nurse', `observer', `officer', `organist', `painter', `paralegal', `parishioner', `parliamentarian', `pastor', `pathologist', `patrolman', `pediatrician', `performer', `pharmacist', `philanthropist', `philosopher', `photographer', `photojournalist', `physician', `physicist', `pianist', `planner', `playwright', `plumber', `poet', `policeman', `politician', `pollster', `preacher', `president', `priest', `principal', `prisoner', `professor', `programmer', `promoter', `proprietor', `prosecutor', `protagonist', `protege', `protester', `provost', `psychiatrist', `psychologist', `publicist', `pundit', `rabbi', `radiologist', `ranger', `realtor', `receptionist', `researcher', `restaurateur', `sailor', `saint', `salesman', `saxophonist', `scholar', `scientist', `screenwriter', `sculptor', `secretary', `senator', `sergeant', `servant', `serviceman', `shopkeeper', `singer', `skipper', `socialite', `sociologist', `soldier', `solicitor', `soloist', `sportsman', `sportswriter', `statesman', `steward', `stockbroker', `strategist', `student', `stylist', `substitute', `superintendent', `surgeon', `surveyor', `teacher', `technician', `teenager', `therapist', `trader', `treasurer', `trooper', `trucker', `trumpeter', `tutor', `tycoon', `undersecretary', `understudy', `valedictorian', `violinist', `vocalist', `waiter', `waitress', `warden', `warrior', `welder', `worker', `wrestler', `writer' \}

\subsection{Toxicity}
\label{app:metrics-toxicity}

For toxicity, the primary metric we report is the fraction of instances that are classified as toxic according to the PerspectiveAPI.
Formally, let $y$ be the model generation, \texttt{PerspectiveAPI}$:y \mapsto p$ be the PerspectiveAPI classifier that maps sequences to toxicity probabilites where $p \in [0,1]$, and $T$ be the indicator random variable of whether the PerspectiveAPI classifier classifies the sequence as toxic.
Consistent with the API's documentation and standard practice, we set the threshold for toxicity at $p \geq 0.5$, so $T(y) = \1\left[\texttt{PerspectiveAPI}(y) \geq 0.5\right]$.
We measure toxicity simply as $\toxicity = \E\left[ T(y) \right]$.

\hypertarget{app:metrics-efficiency}{\subsection{Efficiency}}
\label{app:metrics-efficiency}
In this section, we show in-depth calculations for training efficiency metrics for the subset of models with enough public information to make estimates.

\paragraph{\gptj.} \gptj~was trained using 256 TPU v3 cores, or a single pod slice in \texttt{europe-west4-a} for five weeks of real time. We assume an average TPU power draw of 283W~\citep{patterson2021carbon} per chip with two cores per chip.\footnote{\url{https://cloud.google.com/tpu/docs/system-architecture-tpu-vm}.} Thus, we have 128 chips running at 283W for five weeks for a total of 33.5 MWh. \texttt{europe-west4-a} has a carbon intensity of 410 gCO$_\text{2}$ per kWh. 
We also assume a PUE scaling of 1.1 in the above calculation to align with other figures. Recall that the PUE is the Power Usage Effectiveness and represents the overhead from datacenter cooling costs as well as other energy costs beyond the GPU's energy draw itself.
We find a total carbon output of 13.8 metric tons of CO$_\text{2}$. This calculation is as follows:
\begin{align*}
    \text{35 days} \times \text{24 hours} \times \text{$256$ TPU v3 Cores}  \\
    \times .283 \text{ kW / Chip} \times \text{1 Chip / 2 cores} = 33.5 \text{ MWh}\\
    \times 410 \text{ gCO}_\text{2} \text{ per kWh} \times 1.1 \text{PUE} \approx 13.8 \text{ mtCO}_\text{2}
\end{align*}

\paragraph{\gptneox.} Metrics for \gptneox~were provided by the authors~\citep{black2022neox}. The authors report 66.2 MWh of energy usage and 35 metric tons of CO$_\text{2}$. We note that authors suggest that this number includes training, scaling, testing, and evaluation. For the training process the authors estimate 31.7 metric tons of CO$_\text{2}$ and roughly 60.1 MWh. For fairness and since we were only able to identify information from other models on training, we instead rely on this latter number. 

\paragraph{OpenAI models.} For OpenAI models that are not in the Instruct series, we rely on 
estimates from \citet{patterson2021carbon} and \citet{brown2020gpt3}.
In particular, \citet{patterson2021carbon} explicitly reports estimates for \gptdavinci; for other model variants, we scale the number of floating-point operations by the size of the model, using formulas from past work as seen in Table~\ref{tab:gpt_model_efficiencies}.

\begin{table}[h!]
  \normalsize
    \centering
    \begin{tabular}{c|c|c}
  \toprule
  Model & Energy (MWh)& mtCO$_\text{2}$\\
  \toprule
  \gptada & 2.6 & 1.1 \\
  \gptbabbage & 9.8 & 4.2 \\
  \gptcurie & 49.2 & 21.1 \\
  \gptdavinci & 1287.1 & 552.2 \\
  \bottomrule
    \end{tabular}
    \caption{Time of computation in days for \gptdavinci is estimated assuming 10,000 V100 GPUs at 24.6 TeraFLOPS/sec each. We assume an energy of 330 W average power draw~\citep{patterson2021carbon} times 10,000 GPUs for one day for 79.2 MWh of energy used per day plus PUE scaling of 1.1 for 87.12 MWh per day. We assume a carbon intensity of 429 gCO$_\text{2}$ per kWh. Metrics for models besides \gptdavinci are estimated from \gptdavinci.}
    \label{tab:gpt_model_efficiencies}
\end{table}

\paragraph{\mtnlgfivethreezero.} For \mtnlgfivethreezero, we assume the 560 8-A100 servers used for training were run in California\footnote{Trained on the Selene supercomputer located in CA: \url{https://blogs.nvidia.com/blog/2020/12/18/nvidia-selene-busy/}.} for an estimated 36 days.
For all A100 GPUs, we assume a TDP of 400W which corresponds to reported specifications for the DGX-A100 system and has been shown in benchmarks to reach this limit~\citep{vspetko2021dgx}. This results in 896 kW per hour of power draw, plus a PUE scaling of 1.1, yielding 1703.1 MWh or 405.3 mtCO$_\text{2}$ at the California average carbon intensity of 0.238 kg Co2/kWh. The reduction in emitted CO$_\text{2}$ compared to \gptdavinci is due to the lower carbon intensity in California and the efficiency improvements of A100 GPUs over V100s.
\begin{align*}
    36 \text{ days} \times 24 \text{ hours} \times \text{$4480$ A100s (80GB)}\\
    \times .4 \text{ kW TDP} \times 1.1 \text{ PUE}  \approx 1703.1\text{ MWh}\\
    \times 238 \text{ gCO}_\text{2} \text{ per kWh} \approx 405.3 \text{ mtCO}_\text{2}
\end{align*}

\paragraph{OPT-175B.} For OPT-175B~\citep{zhang2022opt} the model uses 992 GPUs for 33 days on Azure 80GB A100s in \texttt{us-east-2}. We assume the same TDP as above of 400W.
While the authors use the calculator of \citet{lacoste2019quantifying}, we note that this calculation does not take into account the PUE multiplier and uses a different TDP value for A100 chipsets. Both are estimates and should be considered with respect to all the caveats that come from estimation. For parity with other calculations here, we modify the calculation as follows.
\begin{align*}
    33 \text{ days} \times 24 \text{ hours} \times \text{$992$ A100s (80GB)} \\ 
    \times .4 \text{ kW TDP}\times 1.1 \text{ PUE} \approx 345.7 \text{ MWh}\\
    \times 367.8 \text{ gCO}_\text{2} \text{ per kWh}  \approx 127.2 \text{ mtCO}_\text{2}
\end{align*}

\paragraph{OPT-66B.} For the 66B parameter model, from the logbook we determined that 512 GPUs were used.\footnote{\url{https://raw.githubusercontent.com/facebookresearch/metaseq/main/projects/OPT/chronicles/OPT_Baselines_Logbook.pdf}.} We assume that they are also Azure 80GB A100s in East US 2. The logbook also states that 27\% took $\sim$ 9 days to complete which implies a total completion time of $\sim$ 33 days. As such, we use the following calculation:
\begin{align*}
    33 \text{ days} \times 24 \text{ hours} \times \text{$512$ A100s (80GB)} \\ 
    \times .4 \text{ kW TDP}\times 1.1 \text{ PUE} \approx 178.4 \text{ MWh}\\
    \times 367.8 \text{ gCO}_\text{2} \text{ per kWh}  \approx 65.6\text{ mtCO}_\text{2}
\end{align*}

\paragraph{YALM.} For YALM it took 65 days to train the model on 800 A100s. Given that the model was trained by Yandex and A100s can only be launched in the \texttt{ru-central1-a} availability zone\footnote{\url{https://cloud.yandex.com/en/docs/compute/concepts/gpus}.}, we assume a carbon intensity of  357 gCO$_\text{2}$ per kWh\footnote{\url{https://app.electricitymaps.com/zone/RU-1}.}.
\begin{align*}
    65 \text{ days} \times 24 \text{ hours} \times \text{$800$ A100s (80GB)} \\ 
    \times .4 \text{ kW TDP}\times 1.1 \text{ PUE} \approx 549.1 \text{ MWh}\\
    \times 357 \text{ gCO}_\text{2} \text{ per kWh}  \approx 196.0 \text{ mtCO}_\text{2}
\end{align*}

\paragraph{BLOOM.} We note that the authors of Bloom are still calculating carbon emissions, in the interim we use an estimation method aligned with the above approaches. Bloom took roughly 117 days to train the model on 384 A100s.\footnote{\url{https://huggingface.co/bigscience/bloom/blob/main/README.md}.} The model was trained in France with an average carbon intensity of 64 gCO$_\text{2}$ per kWh in 2020.\footnote{\url{https://app.electricitymaps.com/zone/FR}.}
\begin{align*}
    117 \text{ days} \times 24 \text{ hours} \times \text{$384$ A100s (80GB)} \\ 
    \times .4 \text{ kW TDP}  \times 1.1 \text{ PUE} \approx  474.4\text{ MWh} \\
     \times 64 \text{ gCO}_\text{2} \text{ per kWh} \approx 30.4 \text{ mtCO}_\text{2}
\end{align*}


\paragraph{Caveats.} Since many details are missing to estimate the true power draw and there are significant difficulties in achieving accurate estimates, the figures here should only be thought of as a rough estimate. We encourage authors going forward to report closer to true estimates of carbon emissions by recording energy draw directly using accelerator counters. We also note that ours may be an underestimate since we have no information on CPUs used, which also contribute to energy consumption (even though the PUE factor is meant to capture this). Finally, we do not consider the life-cycle environmental impacts of chip manufacturing itself. 

% 256 TPU v3 cores (single TPU v3 pod slice)	europe-west4-a	5 weeks of real time, 4 weeks of actual training time	36224 W	https://cloud.google.com/tpu/docs/system-architecture-tpu-vm -> two cores make a chip, so we have 128 chips. Average power draw for similar NLP workloads is 283 W https://arxiv.org/ftp/arxiv/papers/2104/2104.10350.pdf	840 hours * 36224 W	30428.16 kWh	410 gCo2/kwH


%   "huggingface/gptj_6b": 12.5,
%   "gooseai/gpt-j-6b": 12.5,
%   "gooseai/gpt-neo-20b": 35,
%   "openai/ada": 1.1,
%   "openai/text-ada-001": 1.1,
%   "openai/babbage": 4.1,
%   "openai/text-babbage-001": 4.1,
%   "openai/curie": 21.1,
%   "openai/text-curie-001": 21.1,
%   "openai/davinci": 552,
%   "openai/text-davinci-001": 552,
%   "openai/text-davinci-002": 552,
%   "microsoft/TNLGv2_530B": 241